\section{引言：什么是 Post-Training}

本讲由 Liquid AI 的研究人员（《LLM Engineer's Handbook》作者）主讲，系统介绍了大语言模型（LLM）的 \textbf{Post-Training}（后训练）流程。Post-Training 是预训练之后的所有训练步骤，目标是将一个只会预测下一个 token 的 base model 转化为一个真正有用的 AI assistant。

在预训练阶段，模型在海量原始文本上以 next-token prediction 为目标进行训练，产生的 base model 虽然具备强大的语言建模能力，但并不具备回答问题或遵循指令的能力。Post-Training 正是弥补这一鸿沟的关键环节。

\begin{importantbox}{Post-Training 的两个核心阶段}
\begin{enumerate}
    \item \textbf{Supervised Fine-Tuning (SFT)}：通过指令-回答对（instruction-answer pairs）教会模型遵循指令、以对话形式回答问题。模型在此阶段学习 chat template 结构（system / user / assistant）。
    \item \textbf{Preference Alignment（偏好对齐）}：通过对比 chosen answer（期望回答）与 rejected answer（不期望回答）来调整模型行为，使其产出更符合人类偏好的回答。
\end{enumerate}
\end{importantbox}

\begin{knowledgebox}{Pre-Training vs. Post-Training 的数据规模差异}
预训练使用数万亿（trillions）tokens 的原始文本，数据集规模极为庞大。而 Post-Training 的数据集相对小得多，通常在数千到百万量级，但对数据\textbf{质量}的要求远高于数量。这一差异是理解 Post-Training 方法论的基础。
\end{knowledgebox}

\subsection{本章小结}

Post-Training 是将 base model 转化为实用 assistant 的核心环节，包含 SFT 和 Preference Alignment 两个主要阶段。它的特点是数据规模远小于预训练，但对质量要求极高。

\newpage
